\documentclass[11pt,a4paper]{article}

\usepackage[utf8]{inputenc}
\usepackage[T1]{fontenc}
\usepackage{lmodern}
\usepackage[a4paper,margin=2.4cm]{geometry}
\usepackage{graphicx}
\usepackage{xcolor}
\usepackage{listings}
\usepackage{booktabs}
\usepackage{amsmath}
\usepackage{caption}

% Localisation française manuelle (babel-french indisponible sur cette installation).
\renewcommand{\contentsname}{Table des matières}
\renewcommand{\figurename}{Figure}
\renewcommand{\tablename}{Table}
\frenchspacing

% --- Liens et index : tout en NOIR (hidelinks) ---
\usepackage[hidelinks]{hyperref}
\hypersetup{
  pdftitle={Musee virtuel en Babylon.js -- Rapport de projet ReV},
  pdfauthor={Juan Marcos De Miguel Funes et Julian Rayes},
  colorlinks=false
}

\captionsetup{font=small,labelfont=it,textfont=it}

% --- Style des extraits de code ---
\definecolor{codegray}{gray}{0.95}
\definecolor{codecomment}{rgb}{0.40,0.45,0.40}
\definecolor{codekw}{rgb}{0.15,0.15,0.55}
\definecolor{codestr}{rgb}{0.55,0.20,0.10}
\lstset{
  basicstyle=\ttfamily\footnotesize,
  backgroundcolor=\color{codegray},
  commentstyle=\color{codecomment}\itshape,
  keywordstyle=\color{codekw}\bfseries,
  stringstyle=\color{codestr},
  numberstyle=\tiny\color{gray},
  numbers=left, numbersep=7pt,
  frame=single, rulecolor=\color{gray!40},
  breaklines=true, showstringspaces=false, tabsize=2, captionpos=b,
  language=Java,
  morekeywords={let,const,var,function,class,extends,return,new,this,import,export,from,null,true,false,of,forEach}
}

% --- Emplacement réservé pour une figure (capture à insérer) ---
\newcommand{\figslot}[2][0.70\linewidth]{%
  \begin{figure}[ht]\centering
  \setlength{\fboxsep}{0pt}%
  \fcolorbox{gray!50}{gray!8}{%
    \begin{minipage}[c][3.8cm][c]{#1}\centering\itshape\small
      [\,Emplacement figure --- capture d'écran à insérer\,]
    \end{minipage}}%
  \caption{#2}
  \end{figure}}

\setlength{\parskip}{0.45em}
\setlength{\parindent}{0pt}
\setlength{\emergencystretch}{3em}

\begin{document}

% ============================ PAGE DE TITRE ============================
\begin{titlepage}
\centering
\vspace*{1.5cm}
{\Large\bfseries École Nationale d'Ingénieurs de Brest\par}
\vspace{0.3cm}
{\large Module ReV --- Réalité Virtuelle\par}
\vspace{2.5cm}
\rule{\linewidth}{0.5pt}\\[0.5cm]
{\huge\bfseries Conception d'un musée virtuel\\[0.3cm] en Babylon.js\par}
\vspace{0.4cm}
{\large\itshape Modélisation, animation et simulation de foule\par}
\rule{\linewidth}{0.5pt}\\[2.5cm]

{\large
\textbf{Juan Marcos De Miguel Funes}\\[0.2cm]
\textbf{Julian Rayes}\par}
\vspace{2.5cm}
{\large Encadrant : Éric Maisel\par}
\vfill
{\large Brest --- Juin 2026\par}
\end{titlepage}

% ============================ TABLE DES MATIERES ============================
\clearpage
\tableofcontents
\clearpage

% ============================ 1. INTRODUCTION ============================
\section{Introduction}

Ce rapport décrit la conception d'une \emph{visite de musée virtuel} réalisée
avec le moteur 3D temps réel \textbf{Babylon.js}, dans le cadre du module ReV.
Le bâtiment occupe une emprise de $30\times30$~m~: un hall vitré de $15\times30$~m
et de $10$~m de haut au sud, et au nord trois salles fermées par des portes
coulissantes, surmontées d'une mezzanine accessible par un escalier. Les salles
exposent des tableaux~; le hall et la mezzanine accueillent des statues et un
mobile articulé.

Le projet a été développé en groupe et étalé sur plusieurs séances. Une première
base a été posée --- squelette de l'application (couches \texttt{Visu}/\texttt{Simu}),
système entités/composants, premières primitives, bâtiment initial et tableaux ---
puis progressivement étendue~: portes automatiques, sculptures, décor extérieur,
garde-corps, et enfin une simulation de visiteurs autonomes et de foule guidée. La
collection d'oeuvres est volontairement \emph{éclectique}, mêlant chefs-d'oeuvre
classiques (la Joconde, la Création d'Adam, la Nuit étoilée) et images plus
personnelles~; elle sert de support aux fonctionnalités techniques qui font le
coeur du projet.

L'application utilise les modules ES~; elle se lance en servant le dossier en HTTP
(\texttt{python3 -m http.server}) puis en ouvrant \texttt{index.html}, qui charge
Babylon, instancie la classe \texttt{World} et démarre la boucle de rendu.

\figslot{Vue d'ensemble du musée depuis l'extérieur~: hall vitré, parvis et ciel.}

% ============================ 2. ARCHITECTURE ============================
\section{Architecture logicielle}

Le projet repose sur une architecture \textbf{entités--composants (ECS)} inspirée
du cours \emph{Simuler pour animer}. Une \texttt{Entity}
(\texttt{js/entities/entities.js}) est un conteneur --- nom, position,
\texttt{object3d}, liste de composants --- doté d'une API chaînable
(\texttt{createEntity(...).add(...)}). Deux spécialisations gèrent le mouvement~:
\texttt{Kine} (cinématique) et \texttt{Newton} (forces $\rightarrow$ accélération
$\rightarrow$ vitesse $\rightarrow$ position, par intégration d'Euler).

Le point clé est la \textbf{double passe} de \texttt{Simu.update(dt)}~: chaque
image, on \emph{décide} d'abord (tous les composants calculent leurs forces et
font évoluer les arbres de comportement), puis on \emph{intègre} (le mouvement
résultant est appliqué). Un composant d'IA ne touche donc jamais l'état physique
en cours de calcul.

\begin{lstlisting}[caption={Double passe décision / intégration (js/simu.js).},label=lst:simu]
update(dt) {
  this.entities.forEach((e) => { if (e.execute) e.execute(dt); }); // 1. decision
  this.entities.forEach((e) => { if (e.update)  e.update(dt);  }); // 2. integration
}
\end{lstlisting}

Les composants ne se connaissent pas entre eux~: ils communiquent par un
\emph{blackboard} (\texttt{entity.blackboard}). La trajectoire y écrit le point
visé, le steering le lit, l'arbre de comportement y stocke l'état de pause, le
guide y publie le tableau présenté. Ce découplage a permis d'empiler de nouveaux
comportements sans réécrire l'existant.

% ============================ 3. MODELISATION ============================
\section{Modélisation et matériaux}

\subsection{Structure, ouvertures et collisions}

Les cloisons sont produites par \texttt{js/prims.js}. Les ouvertures (portes,
baies) sont creusées par \textbf{opérations CSG}~: la fonction \texttt{creuser}
soustrait au mur un volume-outil (\texttt{BABYLON.CSG.subtract}). Le mur intérieur
qui sépare le hall des salles est ainsi percé de trois portes, et les baies
vitrées reçoivent un cadre et une vitre semi-transparente. La caméra est une
\texttt{UniversalCamera} en vue première personne, avec capsule de collision,
gravité et déplacements QWERTY~; tous les éléments structurels sont collisionnables.

\figslot{Hall vu de l'intérieur~: baies vitrées et mur à trois portes coulissantes.}

\subsection{Escalier, garde-corps et habillage}

L'escalier (20 marches) double chaque marche~: une \emph{planche visible} fine et
un \emph{bloc de collision invisible} dont la hauteur croît avec l'indice. C'est
ce bloc qui est collisionnable, ce qui évite à la caméra de rester bloquée ou de
sauter de marche en marche --- problème classique des escaliers. Le bord ouvert
de la mezzanine est protégé par un \textbf{garde-corps en verre} collisionnable,
interrompu là où débouche l'escalier pour laisser le passage.

Côté matériaux, le musée combine teintes unies (murs gris à deux variantes,
tapis, piédestaux) et textures (sol intérieur, skybox). Le parvis extérieur est
une \emph{pelouse procédurale}~: une \texttt{DynamicTexture} $16\times16$ remplie
au hasard d'une palette de verts et échantillonnée en mode \texttt{NEAREST}
(rendu « bloc » façon Minecraft), sans aucune dépendance à un fichier image.

% ============================ 4. TABLEAUX ET SCULPTURES ============================
\section{Tableaux et sculptures}

\subsection{Tableaux}

Chaque tableau est un \emph{poster} double-face (deux plans dos à dos partageant
la même texture). Il est \emph{accroché} en devenant \texttt{parent} du mur cible,
puis positionné en \textbf{coordonnées locales}~: un décalage perpendiculaire le
décolle du mur et choisit la face exposée, une rotation l'oriente vers
l'intérieur. Les positions évitent les baies vitrées. Une quinzaine de tableaux
sont répartis dans les salles et sur la mezzanine.

\figslot{Une salle d'exposition~: tableaux accrochés aux murs.}

\subsection{Statues importées et mobiles animés}

Le musée présente trois familles de sculptures.

\textbf{Statue tournante.} La sculpture \emph{La Valse} (\texttt{.glb}) est
importée via \texttt{ImportMesh} et tourne lentement sur un piédestal. La fonction
d'import dispose d'une option \texttt{poserAuSol} qui, après mise à l'échelle,
calcule la boîte englobante pour poser la base au sol et centrer l'empreinte du
modèle sur son axe --- sans quoi elle traversait le sol et tournait de travers.

\figslot{La statue \emph{La Valse}, importée et tournant sur son piédestal.}

\textbf{Maillages statiques.} Trois animaux low-poly (chat, cheval, loup) ornent
le hall d'entrée, alignés face à la porte et rendus collisionnables (obstacles).
Leur orientation est fixée dans le \emph{callback} \texttt{onCharge}, exécuté
\emph{après} le centrage automatique --- la fixer avant aurait décentré le modèle.

\figslot{Les animaux low-poly importés dans le hall, faisant face à l'entrée.}

\textbf{Mobile articulé (Calder).} Au-dessus de la mezzanine pend un mobile
entièrement géométrique, construit \emph{récursivement}~: chaque niveau est une
tige pivotant sur un axe vertical, portant un disque coloré d'un côté et le niveau
inférieur (suspendu par un fil) de l'autre. Le composant \texttt{Mobile} fait
tourner chaque tige à sa propre vitesse, en sens opposé d'un niveau à l'autre~;
le mouvement est donc \emph{articulé} (les parties bougent les unes par rapport aux
autres), contrairement à la rotation d'un bloc rigide de \emph{La Valse}.

\figslot{Le mobile articulé (style Calder) suspendu au-dessus de la mezzanine.}

% ============================ 5. PORTES ============================
\section{Portes coulissantes automatiques}

Les quatre portes (entrée principale et trois portes internes) s'ouvrent à
l'approche d'un agent et se referment à son éloignement, en suivant la directive
du cours (« regarder régulièrement »)~: la proximité est évaluée à \emph{chaque
pas}, pas par un raycast ponctuel. Chaque battant mémorise sa position fermée et
ouverte~; le composant \texttt{AutoDoor} mesure la distance de l'agent le plus
proche (caméra ou visiteur) et interpole en douceur les battants.

\begin{lstlisting}[caption={Ouverture par proximité (js/components/autoDoor.js).},label=lst:autodoor]
execute(dt){
  const open = this.nearestAgentDistance() < this.radius;
  for (const p of this.panels) {
    const targetX = open ? p.openX : p.closedX;
    p.mesh.position.x += (targetX - p.mesh.position.x) * this.speed; // lissage
  }
}
\end{lstlisting}

Le dimensionnement des battants a demandé quelques itérations~: fermés, les
battants laissaient voir l'intérieur par les jointures. La version finale les
surdimensionne uniquement vers leurs bords \emph{exposés} (extérieur et haut), en
laissant le bord intérieur (là où les deux battants se rejoignent) au centre exact,
sans quoi ils se chevauchaient.

\figslot{Une porte coulissante s'ouvrant à l'approche d'un visiteur.}

% ============================ 6. VISITEURS IA ============================
\section{Visiteurs autonomes : steering et arbres de comportement}
\label{sec:visiteurs}

Le coeur dynamique du musée est une population de visiteurs autonomes~: ils
marchent dans une salle, franchissent une porte, s'arrêtent devant une oeuvre,
puis repartent en boucle. Chaque visiteur combine trois mécanismes indépendants
reliés par le blackboard.

\subsection{Steering et trajectoire}

\texttt{Seek} et \texttt{Arrive} (d'après Reynolds, 1987) appliquent une force de
pilotage $\vec{F_s}=k(\vec{v_d}-\vec{v_c})$ vers le point visé~; \texttt{Arrive}
ralentit dans un rayon pour s'arrêter sans osciller, et c'est lui qu'utilisent les
visiteurs. \texttt{Trajectory} enchaîne des points de passage (avec bouclage,
point de départ et sens réglables, pour désynchroniser plusieurs visiteurs sur le
même circuit). Ces points sont calculés à partir des coordonnées réelles des murs
et des portes, pour que chaque segment traverse bien une ouverture. Enfin, comme
la caméra, chaque visiteur possède un \textbf{collider invisible} déplacé par
\texttt{moveWithCollisions}~; le corps visible (peau, vêtements, cheveux, yeux)
en est l'enfant, et Babylon le fait glisser le long des murs.

\subsection{Arbre de comportement}

Le module \texttt{js/bt/} fournit les primitives classiques~: \texttt{Selector}
(« OU »), \texttt{Sequence} (« ET »), décorateurs, \texttt{Condition} et
\texttt{Action}, renvoyant chacune un \texttt{Status} (\texttt{SUCCESS},
\texttt{FAILURE}, \texttt{RUNNING}).

\begin{lstlisting}[caption={Noeuds composites Selector et Sequence (js/bt/behaviorTree.js).},label=lst:bt]
class Selector extends Node {            // "OU" : premier non-FAILURE
  tick(entity, dt) {
    for (const child of this.children) {
      const s = child.tick(entity, dt);
      if (s !== Status.FAILURE) return s;
    }
    return Status.FAILURE;
  }
}
class Sequence extends Node {            // "ET" : premier non-SUCCESS
  tick(entity, dt) {
    for (const child of this.children) {
      const s = child.tick(entity, dt);
      if (s !== Status.SUCCESS) return s;
    }
    return Status.SUCCESS;
  }
}
\end{lstlisting}

L'arbre du visiteur ne gouverne que les \emph{transitions d'état} (marcher
$\leftrightarrow$ en pause)~: une branche « est en pause » décompte un minuteur
puis avance au point suivant, une branche « est arrivé » déclenche la pause. Si
aucune ne se déclenche, le Selector renvoie \texttt{FAILURE} --- c'est-à-dire
« continuer à marcher », ce que le steering fait déjà seul. Choisir un arbre de
comportement plutôt qu'une machine à états le rend \emph{composable}~: ajouter un
état revient à ajouter une branche.

\figslot{Plusieurs visiteurs parcourant le musée et s'arrêtant devant les tableaux.}

% ============================ 7. FOULE ============================
\section{Simulation de foule : guide et groupe}

La dernière étape étend les visiteurs en une \emph{simulation de foule}, en
s'appuyant sur le cours consacré aux comportements collectifs, et toujours
par-dessus l'infrastructure ECS existante.

\subsection{Règles de Reynolds}

Trois composants implémentent le flocking, chacun lisant ses voisins dans un
tableau partagé~: \texttt{Separation} (répulsion $\propto 1/d^2$),
\texttt{Cohesion} (attraction vers le barycentre du groupe) et \texttt{Alignment}
(alignement sur la vitesse moyenne).

\begin{lstlisting}[caption={Force de séparation, répulsion en 1/d2 (js/components/separation.js).},label=lst:sep]
execute(dt){
  const force = new BABYLON.Vector3(0,0,0);
  const p = this.entity.position;
  for (const other of this.group){
    if (other === this.entity || !other.position) continue;
    const away = p.subtract(other.position);     // pointe a l'oppose du voisin
    const dist = away.length();
    if (dist > 0.0001 && dist < this.radius)
      force.addInPlace(away.scale(1.0 / (dist * dist)));
  }
  this.entity.applyForce(force.scale(this.k));
}
\end{lstlisting}

\subsection{Guide et suiveurs}

Un \textbf{guide} (entité plus grande, dorée) suit sa propre trajectoire et
s'arrête devant certains tableaux~; à l'arrêt, son arbre de comportement publie le
tableau présenté dans le blackboard, les points sans arrêt étant simplement
traversés pour mener le groupe d'une salle à l'autre. Quatre \textbf{suiveurs}
visent, via \texttt{FollowGuide}, un point situé \emph{derrière} le guide (et non
le guide lui-même, ce qui les ferait s'empiler), avec un freinage type
\emph{Arrive}. Quand le guide présente une oeuvre, le composant \texttt{Gaze}
tourne chaque suiveur vers le même tableau~; le guide reparti, l'orientation
revient à la direction de marche. Les murs et le parcours du guide confinant déjà
le groupe, aucun mécanisme de rebond n'a été nécessaire. La scène compte au final
un guide, quatre suiveurs et six visiteurs libres (trois en bas, trois sur la
mezzanine).

\figslot{Le guide (doré) présentant un tableau~; le groupe se rassemble et tourne
le regard vers l'oeuvre.}

% ============================ 8. PROBLEMES ============================
\section{Problèmes rencontrés et solutions}
\label{sec:problemes}

\textbf{Visiteurs qui flottaient.} Avec un groupe nombreux, les personnages se
tassaient, se grimpaient dessus et finissaient en l'air, et bloquaient les portes
(2~m). En cause~: leurs colliders entraient en collision \emph{entre eux}, et les
forces de cohésion les poussaient vers le haut sans gravité pour les rabaisser. La
solution a été de \emph{renoncer aux collisions entre visiteurs} (groupes de
collision~: ils ignorent les autres personnes mais heurtent murs et portes) et de
confier l'espacement à la force de séparation de Reynolds, avec deux garde-fous~:
un mode \texttt{planar} (vitesse et position recollées au sol) et un plafond de
vitesse.

\textbf{Le guide traversait les visiteurs.} Conséquence directe~: sans collision
mutuelle, le guide passait à travers le groupe. Plutôt que de rétablir les
collisions, une \emph{répulsion d'espace personnel} (composant \texttt{AvoidGuide},
et logique équivalente dans \texttt{FollowGuide}) écarte les visiteurs du chemin du
guide dans un petit rayon. Le déclenchement utilise la distance 3D (un guide à un
autre étage n'affecte personne) et la poussée reste horizontale.

\textbf{Le cheval, ancien centre des trajectoires.} Devenu obstacle solide, le
cheval occupait le point central vers lequel revenaient toutes les trajectoires~:
les visiteurs s'y agglutinaient et le guide restait coincé. Le point de ralliement
a été déplacé derrière le cheval, hors de son empreinte de collision, comme les
positions de départ du groupe.

D'autres ajustements plus ponctuels (statue traversant le sol, disques du mobile
qui semblaient flotter faute de fils de liaison, dimensionnement des portes) ont
été diagnostiqués à partir de captures d'écran prises pendant les tests.

% ============================ 9. CONCLUSION ============================
\section{Conclusion}

Le projet aboutit à un musée virtuel cohérent sur les plans statique et
dynamique~: un bâtiment complet (hall, salles, mezzanine, escalier, baies,
garde-corps), habillé de matériaux variés, peuplé de tableaux, de statues
importées et d'un mobile articulé, doté de portes réactives, et animé par des
visiteurs autonomes jusqu'à une foule guidée. L'apport principal est
\emph{architectural}~: une base entités--composants soignée a permis d'empiler des
comportements (steering, arbres de comportement, flocking, évitement) par simple
composition, et de corriger les problèmes au fur et à mesure sans jamais réécrire
l'existant.

\end{document}
